% =============================================================================
% Anhang A: Arbeitsablauf und Dateistruktur von EvAX
% =============================================================================
\chapter{Arbeitsablauf und Dateistruktur von EvAX}
\label{ch:anhang-workflow}

Dieser Anhang beschreibt die praktische Funktionsweise von \evax{}
und das Zusammenspiel mit dem ab-initio-Streucode \feff{}.
Er dokumentiert die Struktur der Ein- und Ausgabedateien und soll
den konkreten Umgang mit der Software nachvollziehbar machen.


% ─────────────────────────────────────────────────────────────────────────────
\section{Eingabedaten}

\evax{} benötigt zwei Arten von Eingaben:
experimentelle $\chi(k)$-Daten und ein atomistisches Startmodell.

\subsection{Experimentelle $\chi(k)$-Dateien}

Pro Absorptionskante wird eine zweispaltige ASCII-Datei erwartet.
Die erste Spalte enthält den Wellenzahlvektor $k$ in
\si{\per\Angstrom}, die zweite das EXAFS-Signal $\chi(k)$
(dimensionslos).
Bei einer Dreikanten-Analyse (Co, Ni, Cu) liegen also drei Dateien
vor:
\begin{center}
  \texttt{MEA\_Co\_chi.dat},\quad
  \texttt{MEA\_Ni\_chi.dat},\quad
  \texttt{MEA\_Cu\_chi.dat}
\end{center}

Wie diese Daten aus dem gemessenen Absorptionsspektrum $\mu(E)$
gewonnen werden (Normierung, Hintergrundsubtraktion, $E \to k$-Umrechnung),
ist in Abschnitt~\ref{sec:datenvorverarbeitung} beschrieben.

\subsection{Startmodell (POSCAR-Format)}
\label{sec:anhang-poscar}

Das Strukturmodell wird im POSCAR-Format angegeben, das ursprünglich
aus der Software VASP stammt.
Es definiert die Superzelle durch Gittervektoren und die Positionen
aller Atome in fraktionellen Koordinaten:

\begin{lstlisting}[caption={Aufbau einer POSCAR-Datei
  (\texttt{MEA\_CoNiCu.p1}).}, label=lst:poscar]
CoNiCu MEA fcc a=3.54        <-- Kommentarzeile
1.0                           <-- Skalierungsfaktor
  14.16  0.00   0.00          <-- Gittervektor a1 (Angstrom)
   0.00  14.16  0.00          <-- Gittervektor a2
   0.00   0.00  14.16         <-- Gittervektor a3
Co  Ni  Cu                    <-- Elementnamen
85  85  86                    <-- Anzahl pro Element (= 256)
Direct                        <-- fraktionelle Koordinaten
0.000  0.000  0.000  Cu      <-- x y z Element
0.125  0.125  0.000  Co
0.125  0.000  0.125  Co
...                           <-- 256 Zeilen insgesamt
\end{lstlisting}

Die Gittervektoren bestimmen die Geometrie der Superzelle.
Für eine FCC-Struktur mit Gitterkonstante $a = \SI{3.54}{\Angstrom}$
und $4\times4\times4$-Vervielfachung ergibt sich eine kubische Zelle
mit Kantenlänge $L = 4a = \SI{14.16}{\Angstrom}$ und
$4 \times 4^3 = 256$ Atomen.
Die Elemente werden zufällig, aber stöchiometrisch korrekt auf die
Gitterplätze verteilt.

Die Wahl der Kristallstruktur definiert die Topologie des Gitters.
\evax{} kann Atome verschieben, aber die grundlegende Koordination
nicht ändern -- daher muss die Struktur als Startannahme vorgegeben
werden.
Durch Vergleich der Fitresiduen für verschiedene Startstrukturen
(FCC, BCC, HCP) lässt sich bestimmen, welche die Daten am besten
beschreibt (Tabelle~\ref{tab:startstrukturen-anhang}).

\begin{table}[htbp]
  \centering
  \caption{Verwendete Startstrukturen mit jeweils 256 Atomen.}
  \label{tab:startstrukturen-anhang}
  \small
  \begin{tabular}{llll}
    \toprule
    \textbf{Struktur} & \textbf{Datei} &
    \textbf{Gitterkonstante} & \textbf{Superzelle} \\
    \midrule
    FCC & \texttt{MEA\_CoNiCu.p1}
        & $a = \SI{3.54}{\Angstrom}$
        & $4\times4\times4$ \\
    BCC & \texttt{MEA\_CoNiCu\_bcc.p1}
        & $a = \SI{2.87}{\Angstrom}$
        & $4\times4\times8$ \\
    HCP & \texttt{MEA\_CoNiCu\_hcp.p1}
        & $a = \SI{2.51}{\Angstrom}$, $c = \SI{4.07}{\Angstrom}$
        & $8\times4\times4$ \\
    \bottomrule
  \end{tabular}
\end{table}


% ─────────────────────────────────────────────────────────────────────────────
\section{Die Konfigurationsdatei parameters.dat}
\label{sec:anhang-params}

Alle Einstellungen eines \evax{}-Laufs werden in einer einzigen
Textdatei zusammengefasst.
Sie ist in Blöcke gegliedert, die im Folgenden erläutert werden.

\subsection{Strukturmodell}

\begin{tabular}{@{}p{5cm}p{9.5cm}@{}}
  \texttt{File\_with\_structure} &
    Pfad zur POSCAR-Datei (Startmodell). \\
  \texttt{Supercell} &
    Vervielfachung der Einheitszelle, z.\,B.\ \texttt{1x1x1} wenn
    die Superzelle bereits vollständig aufgebaut ist. \\
  \texttt{Initial\_displacement} &
    Zufällige Auslenkung aller Atome zu Beginn (\si{\Angstrom}), um
    die perfekte Gittersymmetrie zu brechen.
\end{tabular}

\subsection{Optimierung}

\begin{tabular}{@{}p{5cm}p{9.5cm}@{}}
  \texttt{Stop\_after} &
    Zahl der Iterationen in der Screening-Phase, in der die
    Streubasis häufiger aktualisiert wird. \\
  \texttt{Froze\_in} &
    Abbruchkriterium: Der Lauf endet, wenn sich der Residual über
    diese Anzahl an Iterationen nicht mehr verbessert. \\
  \texttt{Number\_of\_states} &
    Zahl der parallelen Strukturkopien im evolutionären Algorithmus
    (typisch: 16). \\
  \texttt{Maximal\_step\_length} &
    Maximale Verschiebung eines Atoms pro Iteration
    (\si{\Angstrom}). \\
  \texttt{Maximal\_displacement} &
    Maximale kumulative Verschiebung eines Atoms gegenüber der
    Startposition (\si{\Angstrom}).
\end{tabular}

\subsection{Experimentelle Daten}

\begin{tabular}{@{}p{5cm}p{9.5cm}@{}}
  \texttt{Number\_of\_spectra} &
    Anzahl der gleichzeitig gefitteten Kanten (hier: 3). \\
  \texttt{File\_with\dots EXAFS\_signal} &
    Pfade zu den $\chi(k)$-Eingabedateien. \\
  \texttt{S02} &
    Amplitudenreduktionsfaktor $S_0^2$ pro Kante.
    Wird mit der Option \texttt{-autoES} automatisch bestimmt. \\
  \texttt{dE0} &
    Energieverschiebung $\Delta E_0$ (\si{\electronvolt})
    zwischen theoretischer und experimenteller Kantenlage.
    Ebenfalls durch \texttt{-autoES} bestimmbar.
\end{tabular}

\subsection{FEFF-Konfiguration}

\begin{tabular}{@{}p{5cm}p{9.5cm}@{}}
  \texttt{FEFF\_path} &
    Pfad zum \feff{}-Binary, das die ab-initio-Streurechnung
    durchführt. \\
  \texttt{N\_legs} &
    Maximale Anzahl der Streuarme pro Pfad.
    2 = nur Einfachstreuung (Absorber $\to$ Nachbar $\to$
    zurück), 4 = einschließlich Dreifach- und
    Vierfachstreuung über Zwischenatome. \\
  \texttt{R\_max\_for\_FEFF} &
    Radius des Clusters um jeden Absorber, für den \feff{}
    Streupfade berechnet (\si{\Angstrom}).
    Größere Werte erfassen mehr Nachbarschalen, erhöhen aber
    die Rechenzeit. \\
  \texttt{Edge} &
    Absorberelemente (z.\,B.\ \texttt{Co Ni Cu}). \\
  \texttt{Edge\_type} &
    Kantentyp (z.\,B.\ \texttt{K K K}). \\
  \texttt{Update\_basis\_every} &
    Häufigkeit der Aktualisierung der FEFF-Streubasis
    (in Iterationen). \\
  \texttt{Clustering\_precision} &
    Schwelle, ab der eine veränderte lokale Umgebung als
    neuer Clustertyp erkannt wird.
\end{tabular}

\subsection{Vergleich Theorie--Experiment}

\begin{tabular}{@{}p{5cm}p{9.5cm}@{}}
  \texttt{Space} &
    Raum, in dem der Vergleich stattfindet:
    \texttt{k} ($k$-Raum), \texttt{R} ($R$-Raum) oder
    \texttt{w} (Wavelet-Raum, zweidimensional in $k$ und $R$). \\
  \texttt{k\_min}, \texttt{k\_max} &
    Grenzen des $k$-Bereichs (\si{\per\Angstrom}). \\
  \texttt{k\_power} &
    Exponent der $k$-Gewichtung ($n$ in $k^n \cdot \chi(k)$). \\
  \texttt{R\_min}, \texttt{R\_max} &
    Grenzen des $R$-Fensters für die Fouriertransformation
    (\si{\Angstrom}).
\end{tabular}


% ─────────────────────────────────────────────────────────────────────────────
\section{Ablauf eines EvAX-Laufs}
\label{sec:anhang-ablauf}

\evax{} wird von der Kommandozeile gestartet:
\begin{lstlisting}[language=bash, numbers=none]
./EvAX.exe parameters.dat -autoES
\end{lstlisting}
Die Option \texttt{-autoES} veranlasst eine Vorstufe, in der $S_0^2$
und $\Delta E_0$ für jede Kante optimiert werden, bevor die
eigentliche Strukturoptimierung beginnt.

\subsection{Initialisierung}

Die Superzelle wird aus der POSCAR-Datei gelesen und 16 Kopien
(\emph{States}) angelegt.
Die Atome werden optional um einen kleinen Betrag
(\texttt{Initial\_displacement}) zufällig ausgelenkt, um die
perfekte Gittersymmetrie zu brechen.
Anschließend führt \evax{} für jeden Absorbertyp eine erste
\feff{}-Rechnung durch, um die Streubasis aufzubauen.

\subsection{Zusammenspiel von EvAX und FEFF}

Das zentrale Element der \evax{}-Methode ist die wiederholte
Berechnung theoretischer EXAFS-Spektren aus einem sich ständig
verändernden Strukturmodell.
Dafür nutzt \evax{} den ab-initio-Code \feff{} (Version~8.5)
\cite{Rehr2010} als externen Rechenkern.
Das Zusammenspiel funktioniert wie folgt:

\begin{enumerate}
  \item \textbf{Clusterextraktion:}
    \evax{} wählt ein Absorberatom (z.\,B.\ ein Co-Atom) und
    extrahiert alle Nachbaratome innerhalb des Radius
    \texttt{R\_max\_for\_FEFF} aus der aktuellen Superzelle.
    Diese lokale Umgebung wird als Cluster in eine \feff{}-Eingabedatei
    (\texttt{feff.inp}) geschrieben.

  \item \textbf{FEFF-Rechnung:}
    \feff{} berechnet für diesen Cluster die Streuphasen und
    -amplituden aller relevanten Pfade (bis zu \texttt{N\_legs}
    Streuarme) unter Verwendung der Muffin-Tin-Approximation für
    das Kristallpotential.
    Die selbstkonsistente Berechnung der Potentiale berücksichtigt
    die tatsächliche chemische Umgebung des Absorbers.
    Das Ergebnis wird in der Binärdatei \texttt{feff.bin} abgelegt.

  \item \textbf{Spektrumberechnung:}
    \evax{} liest die Streupfade aus \texttt{feff.bin} und berechnet
    daraus das theoretische $\chi_\text{calc}(k)$ des betreffenden
    Absorberatoms.
    Über alle Absorberatome gleichen Typs wird gemittelt -- bei
    85~Co-Atomen in der Superzelle ergibt sich das mittlere
    Co-K-EXAFS als Durchschnitt über 85 individuelle Spektren.

  \item \textbf{Clusterrecycling:}
    Eine vollständige \feff{}-Rechnung für alle 256~Atome bei jeder
    Iteration wäre viel zu rechenintensiv.
    Stattdessen baut \evax{} eine Bibliothek bereits berechneter
    Cluster auf.
    Nur wenn sich die lokale Umgebung eines Absorbers über eine
    Schwelle (\texttt{Clustering\_precision}) hinaus verändert hat,
    wird eine neue \feff{}-Rechnung ausgelöst.
    Der Parameter \texttt{Update\_basis\_every} steuert zusätzlich,
    nach wie vielen Iterationen die gesamte Streubasis überprüft wird.
\end{enumerate}

\subsection{Iterationsschleife}

In jeder Iteration wählt \evax{} ein zufälliges Atom und verschiebt
es um einen zufälligen Vektor (maximale Länge:
\texttt{Maximal\_step\_length}).
Aus den neuen Atompositionen wird das theoretische EXAFS-Signal
aller Kanten berechnet und im gewählten Vergleichsraum (typischerweise
Wavelet) mit dem Experiment verglichen.
Der Residual quantifiziert die Abweichung als gewichtete Summe über
alle Kanten:
\begin{equation}
  R = \sum_{i\,\in\,\text{Kanten}} w_i \cdot
      \frac{\|\chi_\text{calc}^{(i)} - \chi_\text{exp}^{(i)}\|}
           {\|\chi_\text{exp}^{(i)}\|}
\end{equation}
Die Gewichte $w_i$ werden automatisch angepasst
(\texttt{adjust\_spectra\_weights}), damit keine Kante den Fit
dominiert.

Ob eine Verschiebung akzeptiert wird, entscheidet das
Metropolis-Kriterium:
\begin{itemize}[nosep]
  \item Sinkt der Residual, wird die Verschiebung akzeptiert.
  \item Steigt er, wird sie mit einer Wahrscheinlichkeit akzeptiert,
        die mit dem Betrag der Verschlechterung abnimmt.
        Dies verhindert, dass der Algorithmus in lokalen Minima
        stecken bleibt.
\end{itemize}

\subsection{Evolutionäre Selektion}

In regelmäßigen Abständen werden die 16~States anhand ihres Residuals
verglichen.
Strukturen mit schlechtem Fit werden durch leicht mutierte Kopien
guter Strukturen ersetzt.
Dieses Selektions-Mutations-Schema entspricht dem Prinzip der
natürlichen Auslese und erhöht die Wahrscheinlichkeit, das globale
Minimum des Residuals zu finden.

\subsection{Konvergenz und Abbruch}

Der Lauf endet, wenn der Residual sich über \texttt{Froze\_in}
aufeinanderfolgende Iterationen nicht verbessert hat.
Zusätzlich kann \texttt{Stop\_after} als Obergrenze für die
Screening-Phase gesetzt werden, nach der die Optimierung mit
seltenerer Basisaktualisierung fortfährt.


% ─────────────────────────────────────────────────────────────────────────────
\section{Dateien im Arbeitsverzeichnis}

Während eines Laufs erzeugt \evax{} eine Reihe von Dateien und
Unterordnern im Arbeitsverzeichnis.
Tabelle~\ref{tab:feff-dateien} erläutert die FEFF-bezogenen Dateien,
Tabelle~\ref{tab:evax-dateien} die sonstigen Zwischendateien.

\subsection{FEFF-bezogene Dateien}

\begin{table}[htbp]
  \centering
  \caption{Dateien und Ordner, die mit der \feff{}-Berechnung
    zusammenhängen.}
  \label{tab:feff-dateien}
  \small
  \begin{tabularx}{\textwidth}{lX}
    \toprule
    \textbf{Datei/Ordner} & \textbf{Bedeutung} \\
    \midrule
    \texttt{feff} &
      Symlink auf das \feff{}-Binary (\texttt{feff85L}).
      Berechnet Streuphasen und -amplituden ab initio unter
      Verwendung der Muffin-Tin-Approximation. \\
    \texttt{run-FEFF.exe} &
      Shell-Skript (Wrapper), das \feff{} im richtigen
      Unterverzeichnis aufruft.
      Typischer Inhalt: \texttt{cd \$1 \&\& ../feff}. \\
    \texttt{feff\_0/}, \texttt{feff\_1/}, \ldots &
      Ein Unterordner pro Absorbertyp (0~=~erste Kante, 1~=~zweite
      usw.).
      \evax{} legt hier die \feff{}-Eingabe ab und liest die
      Ergebnisse aus. \\
    \texttt{feff\_0/feff.inp} &
      Eingabedatei für \feff{}: kartesische Koordinaten des Clusters
      um einen Absorber, Muffin-Tin-Radien, Berechnungsparameter
      (Pfadfilter, $R_\text{max}$, $N_\text{legs}$). \\
    \texttt{feff\_0/feff.bin} &
      Binäre Ausgabe von \feff{}: enthält alle berechneten
      Streupfade mit Amplitude, Phasenverschiebung und effektiver
      Pfadlänge.
      Wird von \evax{} gelesen, um $\chi_\text{calc}(k)$
      zusammenzusetzen. \\
    \texttt{feff\_0/chi.dat} &
      Von \feff{} berechnetes $\chi(k)$ für einen einzelnen
      Cluster. \\
    \texttt{feff\_0/feff.log} &
      Protokoll der Rechnung (Potentialberechnung, Pfadsuche,
      Konvergenz der selbstkonsistenten Potentiale). \\
    \texttt{feff\_0/atoms.dat} &
      Atomliste des Clusters im \feff{}-internen Format. \\
    \bottomrule
  \end{tabularx}
\end{table}

\subsection{Sonstige Dateien}

\begin{table}[htbp]
  \centering
  \caption{Weitere Dateien, die \evax{} während eines Laufs erzeugt.}
  \label{tab:evax-dateien}
  \small
  \begin{tabularx}{\textwidth}{lX}
    \toprule
    \textbf{Datei} & \textbf{Bedeutung} \\
    \midrule
    \texttt{parameters.dat} &
      Zentrale Konfigurationsdatei (vgl.\
      Abschnitt~\ref{sec:anhang-params}).
      Wird bei Verwendung von \texttt{-autoES} um die bestimmten
      $S_0^2$- und $\Delta E_0$-Werte ergänzt. \\
    \texttt{window\_K\_0}, \texttt{\_1}, \texttt{\_2} &
      Hanning-Fensterfunktionen für den $k$-Raum, eine pro Kante. \\
    \texttt{window\_R\_0}, \ldots &
      Fensterfunktionen für den $R$-Raum. \\
    \texttt{window\_W\_0}, \ldots &
      Wavelet-Fensterfunktionen (deutlich größer, da
      zweidimensional in $k$ und $R$). \\
    \texttt{tempfile} &
      Temporäre Zwischendaten der aktuellen Iteration. \\
    \bottomrule
  \end{tabularx}
\end{table}


% ─────────────────────────────────────────────────────────────────────────────
\section{Ausgabedateien}
\label{sec:anhang-output}

Nach Abschluss eines Laufs enthält das Verzeichnis
\texttt{output\_files/} die Ergebnisse.
Pro Kante existiert ein Unterordner
(\texttt{Co/}, \texttt{Ni/}, \texttt{Cu/}).

\begin{table}[htbp]
  \centering
  \caption{Ausgabedateien von \evax{} und ihre Inhalte.}
  \label{tab:output-dateien}
  \small
  \begin{tabularx}{\textwidth}{lX}
    \toprule
    \textbf{Datei} & \textbf{Inhalt} \\
    \midrule
    \texttt{output.dat} &
      Konvergenzverlauf: Jede Zeile enthält die Iterationsnummer
      und den gewichteten Gesamtresidual.
      Dient zur Beurteilung, ob der Lauf konvergiert ist. \\
    \texttt{final.xyz} &
      Kartesische Koordinaten aller 256~Atome in der optimierten
      Struktur.
      Periodische Kopien (für die Visualisierung) sind angefügt;
      die ersten 256~Zeilen sind die realen Atome. \\
    \texttt{initial.xyz} &
      Koordinaten der Ausgangsstruktur vor der Optimierung. \\
    \texttt{restart\textit{N}} &
      Checkpoint nach $N$~Iterationen.
      Enthält den vollständigen Zustand (Parameter, Residuen,
      Atompositionen aller 16~States) und kann als Eingabe für
      einen Folgelauf dienen. \\
    \midrule
    \texttt{Co/EXAFS.dat} &
      Drei Spalten: $k$, $\chi_\text{calc}(k)$,
      $\chi_\text{exp}(k)$.
      Ermöglicht den direkten visuellen und quantitativen
      Vergleich. \\
    \texttt{Co/expFT.dat} &
      Fouriertransformierte: $R$, Re($\tilde\chi$),
      Im($\tilde\chi$). \\
    \texttt{Co/expBFT.dat} &
      Rücktransformierte (Back Fourier Transform) des
      gefilterten Signals. \\
    \texttt{Co/expWT.dat} &
      Wavelet-Transformierte (zweidimensionale Matrix
      in $k$ und $R$). \\
    \texttt{Co/ipolEXAFS.dat} &
      Experimentelles Signal, interpoliert auf das $k$-Gitter
      von \evax{}. \\
    \texttt{Co/RDF\_Co\_Ni.dat} &
      Partielle radiale Verteilungsfunktion $g_\text{Co-Ni}(R)$.
      Pro Absorber-Nachbar-Kombination existiert eine Datei
      (9~Dateien bei drei Elementen). \\
    \bottomrule
  \end{tabularx}
\end{table}


% ─────────────────────────────────────────────────────────────────────────────
\section{Von den Ausgabedateien zu physikalischen Größen}

Die Ausgabedateien liegen als ASCII-Text vor und können direkt
weiterverarbeitet werden.
Aus den partiellen RDF-Dateien (\texttt{RDF\_*}) lassen sich durch
Integration des ersten Peaks die Koordinationszahlen $N_{ij}$ und
aus der Peakposition die mittleren Bindungslängen
$\langle R_{ij} \rangle$ bestimmen.

Der Warren-Cowley-Parameter $\alpha_{ij}$ wird aus der
\texttt{final.xyz}-Struktur berechnet, indem für jedes Absorberatom
die Nachbarn innerhalb eines Cutoff-Radius gezählt und mit der
statistischen Erwartung verglichen werden
(vgl.\ Gl.~\eqref{eq:warren_cowley}).

Der Konvergenzverlauf (\texttt{output.dat}) zeigt, ob der
Residual ein Plateau erreicht hat -- ein Zeichen dafür, dass die
Struktur auskonvergiert ist.
Sinkt der Residual bis zum Ende noch erkennbar, sollten mehr
Iterationen (größeres \texttt{Froze\_in}) gewählt werden.

Die EXAFS- und FT-Dateien erlauben eine visuelle Kontrolle:
Systematische Abweichungen zwischen berechnetem und experimentellem
Spektrum in bestimmten $k$- oder $R$-Bereichen können auf
Modelldefizite hinweisen, die im skalaren Residualwert allein nicht
sichtbar werden.
